Un material alcalí que pot emetre electrons per efecte fotoelèctric presenta una funció de treball d'1,30~eV. Sobre la superfície d'aquest material incideix llum groga amb una longitud d'ona de 500~nm.

\begin{apartat}{1}
Quina freqüència i quina energia tenen els fotons de la llum groga?
\end{apartat}

\begin{apartat}{1}
Quina energia cinètica, en eV, tindran els electrons extrets per aquesta llum groga?
\end{apartat}

\dades{%
  Constant de Planck, $h = 6,63\times 10^{-34}\ \mathrm{J\,s}$.\\
  Massa de l'electró, $m_\mathrm{e} = 9,11\times 10^{-31}\ \mathrm{kg}$.\\
  $1\ \mathrm{eV} = 1,60\times 10^{-19}\ \mathrm{J}$.\\
  Velocitat de la llum, $c = 3,00\times 10^{8}\ \mathrm{m\,s^{-1}}$.}
